\begin{proposition}[The complexified norm and its second derivatives]\hypertarget{tag-0019}{}\label{ari:C.7.11}
\StatementRecord{ari:C.7.11}{proved}{0019}{\SeriesRef{ari:C.7.9}, \SeriesRef{ari:C.7.1}, \SeriesRef{ari:C.7.6}, \SeriesRef{ari:eq:cc-fixedfield}, \SeriesRef{ari:C.7.7}; \citep[Theorem 6.1]{huber2022}}

Use the curve, point and real logarithm of \ref{ari:C.7.9}:
\[
 E:y^2+y=x^3-38x+90,\qquad Q=(4,-2)=\exp_Eu,
 \qquad 0<u<\omega_1/2.
\]
Put $\beta=i\sqrt{19}$, $k=\pi/A$, $s_2=-2$, and
\[
 \theta(z)=\sigma(z)e^{-s_2z^2/2},\qquad
 D(z)=\zeta(z)-s_2z,\qquad
 \mathcal F(z,w)=\sqrt{19}\,\theta(z)\theta(w)e^{-kzw},
 \qquad H=\mathcal F^2.
\]
Then $\mathcal F(z,\bar z)=N(z)^2$, and $\mathcal F$ is an
entire function of two complex variables of order at most two.
It is periodic under $(z,w)\mapsto(z+\omega,w+\bar\omega)$,
$\omega\in\Lambda$. The subgroup
\[
 \Gamma=\{(\alpha u+\omega,\bar\alpha u+\bar\omega):
                   \alpha\in\Z[\beta],\ \omega\in\Lambda\}
 \subset\C^2
\]
has rank four over $\Z$ and lies in the real plane
$\{(z,\bar z):z\in\C\}$. There is one number field $K$ such
that, with $H_0=H(u,u)=N(Q)^4$, every value of $H$ on
$\Gamma$ belongs to $K[H_0]$. More precisely, the nonzero values
are algebraic multiples of the powers $H_0^{\alpha\bar\alpha}$,
as specified in \eqref{ari:eq:cc-fixedfield}.

For every nonzero $\alpha\in\Z[\beta]$ there is an algebraic
number $a_\alpha$ such that
\[
 D(\alpha u)=\bar\alpha D(u)+a_\alpha.
\]
At $p_\alpha=(\alpha u,\bar\alpha u)$ the algebraic constant
direction $V_\alpha=\alpha\partial_z-\bar\alpha\partial_w$
satisfies
\begin{align*}
 V_\alpha\log\mathcal F(p_\alpha)
    &=\alpha a_\alpha-\bar\alpha\bar a_\alpha
                         \in\overline{\Q},\\
 V_\alpha^2\log\mathcal F(p_\alpha)
    &=-\alpha^2\{\wp(\alpha u)+s_2\}
      -\bar\alpha^2\{\wp(\bar\alpha u)+s_2\}
      +2\alpha\bar\alpha k.
\end{align*}
The number $k$ is transcendental. For every nonzero $\alpha$,
the exact normalized second derivative satisfies
\[
 \overline\Q\!\left(\frac{V_\alpha^2H(p_\alpha)}{H(p_\alpha)}\right)
      =\overline\Q(k).
\] In particular,
\[
 (\partial_z-\partial_w)H(u,u)=0,\qquad
 \frac{(\partial_z-\partial_w)^2H(u,u)}{H(u,u)}=4k-8.
\]
Thus the normalized first derivative is algebraic and the
normalized second derivative generates $\overline\Q(k)$. The rank-four sampling set and growth order give the numerical
comparison with \eqref{ari:eq:cc-rankcriterion} proved below.
\end{proposition}



\begin{proof}
\begin{enumerate}[label=\arabic*.,leftmargin=*]
\item The real model and its algebraic differential give real sigma
coefficients and a conjugation-stable lattice. In particular
$\overline{\theta(z)}=\theta(\bar z)$. The discriminant is
$-19^3$, so \ref{ari:C.7.1} gives the asserted diagonal identity.
The sigma function and the exponential of the quadratic
polynomial $-kzw$ have order at most two. From \ref{ari:C.7.6},
\[
 \theta(z+\omega)=\epsilon(\omega)
         e^{k\bar\omega(z+\omega/2)}\theta(z),
 \qquad \epsilon(\bar\omega)=\epsilon(\omega)\in\{1,-1\}.
\]
Multiplying this formula and its conjugate version, the two
signs multiply to one and the exponent cancels the change of
$-kzw$. This proves periodicity on all of $\C^2$.

\item If $\alpha u+\omega=0$, then $[\alpha]Q=O$.
Nontorsion of $Q$ and the fact that a nonzero elliptic
endomorphism is an isogeny imply $\alpha=0$, and then
$\omega=0$. Thus the rank-two module $\Z[\beta]u$ and the
rank-two period lattice are independent over $\Z$.
This proves the rank assertion.

\item For $d=\alpha\bar\alpha$, the squared CM law \ref{ari:C.7.6} gives
the exact fixed-field identity
\begin{equation}\hypertarget{tag-002Q}{}\label{ari:eq:cc-fixedfield}
 H(\alpha u,\bar\alpha u)
   =19^{1-d}F_\alpha(Q)\overline{F_\alpha(Q)}H(u,u)^d.
\end{equation}
All the endomorphisms and their normalized rational functions
$F_\alpha$ are defined over a single number field: choose
a field containing the coefficients of the finitely many
generators of the CM order, and close it under conjugation.
The uniqueness of the Laurent normalization in \ref{ari:C.7.6} gives
the same field of definition for every $F_\alpha$.
Writing that fixed number field as $K$, formula
\eqref{ari:eq:cc-fixedfield} places all values in $K[H_0]$.
The value for $\alpha=0$ is zero, and period shifts do not
change it. Squaring was included because algebraicity of the
individual positive square roots alone would not establish
their membership in one fixed number field.

\item Logarithmically differentiating the squared CM law at $u$
gives
\[
 a_\alpha=\frac{1}{2\alpha}
     \frac{(d/dz)F_\alpha(\exp_Ez)|_{z=u}}{F_\alpha(Q)}
                  \in\overline{\Q}.
\]
The denominator is nonzero since $Q$ is nontorsion. The
derivation $d/dz$ is the algebraic invariant derivation on
$E$, so the displayed derivative has algebraic value.
Since $D(u)$ is real, conjugation gives the corresponding
formula at $\bar\alpha u$. On a neighborhood where
$\mathcal F\ne0$, its logarithmic gradient and Hessian are
\label{layout:I:3}
\[
\begin{gathered}
(\partial_z\log\mathcal F,\partial_w\log\mathcal F)
       =(D(z)-kw,D(w)-kz),\qquad
 \\
\operatorname{Hess}\log\mathcal F=
 \begin{pmatrix}-\wp(z)-s_2&-k\\-k&-\wp(w)-s_2\end{pmatrix}.
\end{gathered}
\]
Evaluation on $(\alpha,-\bar\alpha)$ proves both identities
in the statement. In particular the Gaussian cancels in the
first derivative and reappears in the second derivative.

\item To prove $k$ transcendental, recall the real period formula
$\eta_1=(s_2+k)\omega_1$. If $k$ were algebraic, the
nonzero period vector $(\omega_1,\eta_1)$ of the algebraic
universal vector extension of $E$ would lie on the algebraic
line of slope $s_2+k$. Dividing this period by a sufficiently
large integer supplies a nonzero algebraic torsion exponential.
The analytic subgroup theorem
\citep[Theorem~6.1, arXiv v8, p.~43]{huber2022} would give a one-dimensional
algebraic subgroup with that Lie line. Its projection to $E$ is nonzero, so it splits the vector
extension after an isogeny. The nonsplitting argument in
\cref{ari:C.7.7} excludes this: the coordinates here are
exactly $(z,w)$ with periods $(\omega,\eta(\omega))$ and
algebraic charts $t=w-\zeta(z)$,
$t_O=t-\wp'(z)/(2\wp(z))$.

\item Write the two logarithmic derivatives as
$V_\alpha\log\mathcal F=B_\alpha$ and
$V_\alpha^2\log\mathcal F=A_\alpha+2dk$, where
$A_\alpha,B_\alpha$ are algebraic. Since $H=\mathcal F^2$,
\[
 \frac{V_\alpha^2H(p_\alpha)}{H(p_\alpha)}
       =4B_\alpha^2+2A_\alpha+4dk.
\]
Here $d>0$, so the displayed affine expression both belongs to
$\overline\Q(k)$ and recovers $k$ by subtraction of its algebraic
constant and division by $4d$. This proves the asserted field
equality; $H(p_\alpha)$ is nonzero because $Q$ is nontorsion.
At $\alpha=1$, $a_1=0$, $\wp(u)=4$ and $s_2=-2$,
which give the last displayed formula in the statement.

\item \textbf{Differential and growth criteria.}
The differential criterion of
\citet[Theorems~3.6--3.7, author proof, pp.~21--22]{waldschmidt2009auxiliary}
requires a finite algebraically independent family whose
partial derivatives belong to the polynomial ring generated
by that family over one number field. Algebraic values of
the full family then give algebraic values of all iterated
partial derivatives. The calculation above gives the exact value-field obstruction:
the field containing $H(p_\alpha)$ and its second directional
derivative contains $k$ as soon as the nonzero value is inverted.
In particular its normalized second jet is transcendental. The issue persists
although the first cancelling direction is algebraic; the
direction also varies with $\alpha$.

\item Theorem~1 of
\cite[p.~24]{waldschmidt1975dimension} uses $d$ meromorphic functions on $\C^n$, growth and
arithmetic-size exponents $\rho_1,\ldots,\rho_d$, and sampling
sets of cardinality $N^{\ell-o(1)}$ in a ball of radius $O(N)$,
with separation at least $N^{-\lambda-o(1)}$. Its strict numerical
requirement is
\begin{equation}\hypertarget{tag-002X}{}\label{ari:eq:cc-rankcriterion}
 (d-1)\ell>\sum_{j=1}^d\rho_j
                     +2d(n-1)(\lambda+1).
\end{equation}
The theorem also requires all function values to belong to
one number field and bounds on denominators in the entire
representations of the meromorphic functions. 

\item The natural independent triple here is
$\wp(z),\wp(w),H(z,w)$, all of order at most two.
The first two are algebraically independent. If $H$ satisfied
a polynomial over $\C(\wp(z),\wp(w))$, fix a generic pair
$(z,w)$ and translate only $z$ by $m\omega$, with
$0\ne\omega\in\Lambda$. The coefficients of the polynomial
remain constant, whereas the absolute values of its proposed
roots $H(z+m\omega,w)$ grow as
$\exp(k|\omega|^2m^2+O(|m|))$. This contradicts the finite
set of roots of a nonzero polynomial. Thus the triple is
independent.

\item The growth-order lower envelope is $\rho_j=2$. Substitution of
this envelope and $n=2$, $d=3$, $\ell=4$ into
\eqref{ari:eq:cc-rankcriterion} gives the required inequality
\[
 8>12+6\lambda,
\]
which fails even for $\lambda=0$. Moreover, the $\asymp N^4$
points obtained by bounding
the four integer coefficients lie in a ball of radius $O(N)$
in a real plane. Packing separated discs in that plane gives
$N^4\ll N^{2+2\lambda+o(1)}$, hence $\lambda\ge1$.
Removing the rank-two set of poles does not change this count.
Larger growth exponents only worsen the inequality. Thus the stated rank and growth envelopes do not satisfy
\eqref{ari:eq:cc-rankcriterion}. The norm question is
recorded in \cref{ari:C.7.5}.
\end{enumerate}
\end{proof}